\documentclass[11pt,a4paper]{article}
\usepackage[margin=2.2cm]{geometry}
\usepackage{booktabs}
\usepackage{graphicx}
\usepackage{amsmath}
\usepackage{float}
\usepackage[hidelinks]{hyperref}

\title{Assignment 10: Agentic Systems -- Deep Research Assistant}
\author{}
\date{}

\begin{document}
\maketitle
\vspace{-1.5cm}

\section{System Overview}

We built a deep-research agent that takes a complex research query, decomposes it into sub-queries, searches the web for each, validates results, and synthesizes a markdown report. The system uses DuckDuckGo for web search, OpenRouter's free API (Nvidia Nemotron 120B with model fallback) for LLM calls, and implements three architectural variants.

The pipeline has four stages: (1) \textbf{Decompose}: LLM breaks the query into 3--5 focused sub-queries; (2) \textbf{Search}: web search for each sub-query with result validation; (3) \textbf{Answer}: LLM answers each sub-query from its search results; (4) \textbf{Synthesize}: LLM merges all answers into a coherent markdown report.

\section{Three Architecture Variants}

\textbf{Sequential.} One agent processes everything step by step: decompose, then search and answer each sub-query in order. Simple and deterministic, but slow.

\textbf{Parallel.} Decompose first, then dispatch all sub-query searches concurrently via a thread pool. Assumes all sub-queries are independent.

\textbf{Hierarchical.} A planner agent decomposes the query with explicit dependency analysis, labelling each sub-query as independent or dependent. Independent queries run in parallel (phase 1), dependent queries run sequentially using prior answers as context (phase 2).

\begin{table}[H]
\centering
\begin{tabular}{lccc}
\toprule
Architecture & Time (s) & Report Length (chars) & Sub-queries \\
\midrule
Sequential & 260.8 & 13,852 & 5 \\
Parallel & 128.7 & 12,194 & 4 \\
Hierarchical & 290.4 & 12,866 & 5 \\
\bottomrule
\end{tabular}
\caption{Architecture comparison on query about DeepSeek-R1 vs OpenAI o3.}
\end{table}

Parallel was 2$\times$ faster than sequential because all web searches ran concurrently. Hierarchical was slowest due to the extra planning LLM call and timeout on one sub-query. All three produced substantive reports with real benchmark numbers (e.g.\ DeepSeek-R1: 97.3\% on MATH-500, OpenAI o3: 96.7\% on AIME 2024).

\section{Key Design Decisions}

\textbf{Search validation.} Raw web results are filtered by snippet length ($>$30 chars), keyword overlap with the query, and exclusion of non-content pages. Results are ranked by a relevance score. All queries returned 4--5 validated results, confirming the filter is not too aggressive.

\textbf{Query decomposition.} The Nemotron 120B model consistently produced well-structured JSON sub-queries. For the scaling laws query it generated: (1) DeepSeek-R1 performance metrics, (2) OpenAI o3 performance metrics, (3) test-time vs pre-training scaling comparison, (4) compute costs, (5) cost-effectiveness comparison. This is a sensible decomposition that isolates independent facts before attempting comparisons.

\textbf{Model fallback routing.} OpenRouter's \texttt{models} + \texttt{route: "fallback"} feature automatically tries alternative free models (Step 3.5 Flash, MiniMax M2.5) when the primary model times out, which happened twice during our experiments.

\section{Reflections}

\textbf{Decomposition quality matters most.} The parallel architecture decomposed the first query into 4 sub-queries (vs.\ 5 for sequential/hierarchical), yet produced a comparably informative report. Fewer, more focused sub-queries can be more effective than exhaustive decomposition.

\textbf{Dependency detection was unnecessary for these queries.} The hierarchical planner labelled all 5 sub-queries as independent for the scaling laws query. This is correct---comparing models A and B requires facts about each independently before comparison. For causal reasoning queries (``Why did X cause Y?''), dependencies would matter more.

\textbf{Web search limitations.} Some sub-queries returned results that lacked quantitative data (e.g.\ exact FLOPs for OpenAI o3 training). The agent correctly noted these gaps in its synthesis rather than hallucinating numbers. This is a strength of the RAG-like grounding approach.

\textbf{Parallel is the sweet spot.} Sequential offers no practical advantage (prior results don't inform subsequent searches for comparison queries). Hierarchical adds complexity and latency without benefit when sub-queries are independent. Parallel gives the best time/quality trade-off.

\textbf{Potential improvements:} (1) Iterative refinement---let the synthesizer identify gaps and trigger additional searches; (2) Multi-source verification---cross-check facts across independent sources; (3) Adaptive architecture---detect query type (comparison vs.\ causal) and choose sequential/parallel accordingly; (4) Use OpenRouter's \texttt{web} plugin as an additional search source alongside DuckDuckGo.

\end{document}
